\input{1-preamble}

\author{}

\begin{document}

%TC:ignore
\maketitle

                 
\begin{abstract}
In an asymmetric alliance like NATO, competing policy signals from allied states are consequential for the public's foreign policy attitudes. Three years since the Russian invasion, the initial rally of public support for Ukraine has become a fault line in the transatlantic alliance. Whereas NATO's most powerful member signaled that abandoning an embattled ally was a permissible policy option, European governments have largely maintained their commitment to support Ukraine with military aid. US President Trump’s threats to suspend military aid to Ukraine created a rare window for studying the effects of competing elite cues on transatlantic public opinion. Do competing policy signals from home and abroad strengthen solidarity attitudes in alliances, or do they permit dissent among the public? I fielded a preregistered survey experiment (n = 3029) in Germany and the UK just before the Trump Administration officially halted the flow of military aid to Ukraine, using vignette treatments that anticipated the policy signals eventually made by both countries. The results show that when the domestic government's pro-aid policy is juxtaposed with a countervailing signal from the US, it creates a permissive dissensus. Rather than strengthening solidarity, competing signals from within the alliance caused European citizens to not only become less supportive of increasing military aid for Ukraine, but also less likely to make moral arguments about solidarity.
\end{abstract}

\textbf{Keywords}: Public opinion, alliances, NATO, survey experiment, military aid

%TC:endignore
\newpage

\section{Introduction}
Three years after Russia's full-scale invasion, the transatlantic coalition of the willing has fractured. The 2024 election of Donald Trump and the pivot of US foreign policy away from its transatlantic security commitments has introduced a new source of uncertainty to the future of Western support for Ukraine. Whereas NATO's most powerful member signaled that abandoning an embattled ally was a permissible and legitimate policy option, European governments have largely maintained their commitment to support Ukraine with military aid. 

This particular rift in transatlantic foreign policy created an opportunity to study how competing policy signals from allied states shape citizens' foreign policy attitudes. What happens to public opinion towards military aid for Ukraine when citizens are faced with competing policy signals from home and abroad? 

The consequences of elite discord within an alliance like NATO are particularly salient, especially when the stakes of war are high. NATO's coalition of the willing and the coordination of military aid to Ukraine demonstrates that countries facing heightened military threats can seek security and demonstrate cohesion through collective action. Public opinion research has shown the importance of elite consensus for maintaining public support for the government's security and defense policy. \autocite{berinsky_assuming_2007, groeling_crossing_2008, baum_relationships_2008, baum_shot_2009} When elites send consistent signals, the public tends to follow.\autocite{mueller_war_1973} 

When elite signals diverge, the picture becomes more complicated. Within the domestic sphere, we know that competing cues can cause opinions to polarize along partisan lines. \autocite{berinsky_assuming_2007, druckman_how_2013} However, little is known about how competing cues \emph{across allied states} affect public opinion. NATO's asymmetric structure makes this question especially pressing, given the importance of achieving and communicating consensus on alliance commitments. Do citizens become more or less willing to support military aid for Ukraine when the US withdraws from its commitment and disrupts the consensus? Furthermore, Russian attempts to undermine NATO cohesion and amplify perceptions of elite discord within the alliance makes the relationship between consensus and public support for war all the more interesting to study. 

%Past research suggests that the public punishes leaders who break alliance commitments, especially if it means abandoning an ally that has been unfairly attacked.\autocite{chiba_careful_2015, tomz_military_2021, levin_art_2022} In an experimental setting, alliances have been shown to increase support for intervention both by inciting feelings of loyalty and concerns for one's international and military reputation. However, not all alliance commitments are created equally, and competing narratives about the nature of alliance commitments may cause an initial rally effect to wane. 

This article addresses these gaps by leveraging a preregistered survey experiment (n = 3029)\footnote{A pre-analysis plan was preregistered on the OSF prior to accessing the data, under the following DOI: https://doi.org/10.17605/OSF.IO/CN4V8. The study plan was approved by Aarhus University’s Research Ethics Committee on June 28th, 2024.} in Germany and the United Kingdom, isolating the effects of competing elite policy signals on support for military aid to Ukraine. The survey was fielded just before the Trump Administration officially halted the flow of military aid to Ukraine, using vignette treatments that anticipated the actual policy signals eventually made by Trump and the national governments of both countries. 

This study offers a number of important contributions to the literature on alliances and public support for war. By modeling real-world cue competition just before a major policy shift, I find that disagreement between NATO allies produces a permissive dissensus. Competing signals between NATO allies communicates to the public that the consensus behind supporting Ukraine has fractured, which has a direct impact on public support for military aid. Rather than strengthening solidarity, visible disagreement between Ukraine's principal backers caused European citizens to become less supportive of increasing military aid. Despite attempts by European leaders to reassure Ukraine, it does not appear that domestic governments can counter the negative effect of consensus loss and maintain public support for military aid among their constituents. These findings challenge two assumptions in the literature. First, that the public categorically opposes violating commitments to shared allies,\autocite{chiba_careful_2015, tomz_military_2021, levin_art_2022} and second, that positive arguments can overcome the negative effects of political disagreement about external commitments.\autocite{chu_commanding_2021, reiter_manipulating_2024} 

Furthermore, the study makes use of text analysis to understand how voters justify attitudes towards military aid in their own words. Competing signals caused European citizens to become less likely to make moral arguments about solidarity and invited the public to express grievances about the tradeoffs of sending military aid. The findings highlight the conditional nature of solidarity and alliance commitments, indicating that powerful foreign leaders can signal to the rest of the alliance that abandoning an embattled ally is permissible. In sum, the results underscore the critical role of elite signaling in maintaining credible commitments in times of uncertainty, and emphasize the tenuousness of public support for collective action when alliance cohesion is called into question.

The following section outlines the theory and expectations underlying the study's three pre-registered hypotheses. The research design, methods, and results of the pre-registered survey experiment are then presented and discussed. Finally, I present the results of the exploratory text analysis before summarizing the study's key findings in the conclusion.

\section{How Do Policy Signals Influence Public Support for War?}
\subsection{The Domestic Costs of War}
Public support for war has been shown to fluctuate in response to elite cues and changing circumstances in international security. \autocite{chu_commanding_2021, alley_elite_2023, jentleson_still_1998, gelpi_performing_2010, gelpi_success_2005,  jakobbarrett_how_2026} The high quantity of information about domestic politics as well as international relations about can be overwhelming for individuals. Therefore, the public tends to rely on policy signals from trusted political leaders as heuristic devices that guide opinion formation \autocite{mondak_public_1993, zaller_nature_2006, berinsky_assuming_2007, baum_relationships_2008}. 

Numerous studies have demonstrated the importance of domestic elite consensus for maintaining public support for war among democratic audiences. \autocite{berinsky_assuming_2007, groeling_crossing_2008, baum_relationships_2008, baum_shot_2009} \textcite{berinsky_assuming_2007} argues that when the government presents a unified front in support of a military intervention, public opinion should follow the mainstream pattern, generating an observable rally-effect. In the absence of a unified elite consensus, the scholarship on elite cues predicts that opinions should polarize along partisan lines. \autocite{berinsky_assuming_2007, druckman_how_2013}  Importantly, Berinsky's theories were applied to American public support for active deployments of US troops, such as in World War II and the American invasion of Iraq in 2003. There may be reason to expect different dynamics in the case of European public support for military aid, for two specific reasons. 

First, whereas Berinsky studied rally effects within the US two-party system, both Germany and the UK are multiparty democracies, where a unified elite consensus across multiple parties is more difficult to achieve and hence rarer to observe. We currently lack a clear empirical expectation from the literature about the necessary conditions for rally effects and polarization patterns to emerge in multiparty systems. Consensus among governing parties or mainstream parties may be sufficient to sustain significant public support for war, even if some voters are swayed by ideologically coherent cues from opposition parties.\autocite{brader_which_2013} At the time of fielding this study, a \textit{mainstream}, cross-party consensus in the UK presented domestic audiences with a unified front in support of military aid for Ukraine. In Germany, the pattern was slightly different. At the time of fielding, Germany’s government had fallen amid a disagreement between the FDP and its coalition partners about the federal debt brake. During the 2025 campaign, the strongest voices against military aid came from the AfD and the newly-formed BSW, both of which argued for a much softer approach towards Moscow. However, given that mainstream parties had publicly refused to form a coalition with the AfD, the German government maintained a \textit{governing} consensus in favor of military aid for Ukraine.\footnote{The specificity of the political climate and partisan differences in support for military aid are discussed further in the Appendix.} 

%What is the takeaway, and why is Berinsky maybe not the best theory?
Second, support for military aid constitutes a different form of support for war that falls short of deploying troops to Ukraine. For this reason, I argue that Berinsky's theory of elite consensus may be insufficient to explain public support for military aid. By sending military aid, European democracies are indirectly involved in Ukraine's defensive war against Russia, avoiding the same threats of wartime casualties or loss of territorial sovereignty. While events-response theories of public support for war focus on how informational cues such as rising casualties and the likelihood of success affect public opinion,\autocite{gelpi_paying_2009, gelpi_success_2005, gelpi_iraq_2007} I argue that \emph{public support for military aid} is more likely to follow the patterns that scholars have identified in public opinion towards \textit{defense spending} and \textit{foreign aid}. When the public is informed about plans to increase military aid for Ukraine, they are more likely to interpret this as an indication of more national spending. 
%Info about military aid policy is indicative of more spending, not more casualties!

In spite of mainstream political support for Ukraine at the elite level, there may be a thermostatic effect at the population level when the costs of war feel too much to bear. \textcite{wlezien_public_1995} argues that the public's preferences towards defense and defense spending change systemically in response to information about past policy decisions \autocite{wlezien_public_1995}. In other words, the public acts like a thermostat; when the public receives retrospective information about an increase to defense spending, public preferences for defense spending tend to move in the opposite direction. Importantly, this effect appears to be reactive, and neither assumes nor excludes that opinion change towards defense-spending can be deliberative among individuals with higher political knowledge. \textcite{scotto_we_2017} have revealed a similar effect in relation to public opinion towards foreign aid spending. In general, citizens in high-income countries tend to be skeptical towards foreign aid spending, often assuming the amount of government spending on foreign aid to be higher than it actually is.

The thermostatic model also suggests that when the public is informed about future defense spending policies, they may prospectively associate this policy signal with subsequent declines in public spending. In the real world, individuals are likely to be sensitive to the risks of investing more in \emph{guns than in butter} \autocite{powell_guns_1993, bove_political_2016, carrubba_decision_2004}. In 2023, \textcite{thomson_european_2023} found that Europe's cost-of-living crisis did not seem to affect public attitudes towards sanctions and other indicators of support for Ukraine. However, I argue that thermostatic effects may also fluctuate over time in response to longer term trends in defense spending. Europe's coalition of the willing has been supporting Ukraine with military aid since the invasion began in 2022. With this in mind, the public may react negatively towards information that their government plans to commit even more resources towards military aid for another country. For example, what is seen as \textit{yet another} proposal to increase spending on military aid and defense is more likely to generate concerns about the downstream effects on funding for education and healthcare \autocite{mintz_electoral_1988}. 

Taken together, the insights from the literature offer a mixed picture on whether domestic governments should be able to garner public support for increasing military aid, three years into the war. The clarity of elite consensus in Germany and the UK is much weaker than in \textcite{berinsky_assuming_2007}'s models of US support for war, making it less likely to observe a rally effect. Regarding the costs of war, the public is also much more likely to think about military aid in terms of material tradeoffs, such as the economic and social costs at home. Insofar as the thermostatic hypothesis applies to preferences for military aid, there is little reason to expect that a signal from the domestic government will motivate a fatigued public to increase the amount of weapons deliveries to Ukraine. 

Therefore, I argue that a signal from the domestic government about increasing military aid will not generate an observable rally effect at the public level: 
\paragraph{}
\textit{H1: The domestic policy signal will not increase support for raising the amount of national military aid for Ukraine.}

\subsection{The Cost of Backing Down}
The literature on alliances and public support for war argues that when citizens perceive allied nations as being part of an \textquote{ingroup}, concerns for fairness and loyalty drive public support for war \autocite{tomz_military_2021, tomz_how_2023}. The public shies away from acting disloyally towards foreign allies, and punishes leaders who break alliance commitments with lower approval ratings \autocite{chiba_careful_2015, tomz_military_2021, levin_art_2022} Meanwhile, leaders who retract support for an ally in times of conflict (especially after committing to help in the first place) may pay the price in the form of public disapproval \autocite{tomz_military_2021}. However, these loyalty mechanisms may only operate with respect to formal alliance members, meaning that external partners seeking closer alignment, such as Ukraine with NATO, may not benefit from the same loyalty principles.

Leaders can also avoid public opinion penalties by offering compelling justifications for why backing down is the right decision. \textcite{levendusky_when_2012} show that when the president uses changing circumstances and new information to frame why backing down is in the nation's best interest, the public becomes less likely to express disapproval. I argue that these framing effects are likely to function similarly not just in escalation scenarios, but also when a leader seeks to justify backing down from alliance commitments. However, we lack research that tests this proposition in the context of international alliances, which also engages factors outside of the domestic political sphere. While actions by domestic leaders are the most familiar and salient to public audiences, I argue that highly salient decisions by foreign leaders will also be subject to public opinion penalties from allied audiences. This follows an acknowledgment in the literature that leaders face reputational costs at home as well as abroad. \autocite{levy_backing_2015, agadjanian_has_2020} 

However, the literature still lacks an empirical test of this assumption as it relates to informal commitments. How does the public feel about backing out from informal commitments to countries that are not members of the alliance itself? If public audiences in NATO countries feel a strong sense of ingroup loyalty to Ukraine, the literature provides evidence to suggest that the public will disapprove of abandoning Ukraine. After all, the North Atlantic alliance and its coalition of the willing have publicly (albeit informally) committed to support Ukraine in its defensive war against the Russian invasion. Nevertheless, policy signals from foreign leaders can be consequential for public attitudes towards defense and security, especially if signals from abroad constitute credible threats for alliance partners.\autocite{blankenship_price_2021, blankenship_threats_2024} Recent research suggests that when the US signals a desire to withdraw completely from its alliance commitments, European voters develop a stronger preference for security autonomy. \autocite{jakobbarrett_how_2026} These findings imply that under the pressure of a dominant ally withdrawing, the public may prefer to turn inwards and minimize security burdens. 

The following hypothesis seeks to test whether signals of withdrawal from the alliance's most powerful state persuades citizens to break commitments to a country that lies beyond the formal boundaries of the alliance. Research has documented a high degree of public support for Ukraine in many European countries \autocite{thomson_european_2023, mader_increased_2024}. Furthermore, the predominance of pro-Ukraine narratives among governing parties in Germany and the UK suggests that Trump's plans to cut off military aid for Ukraine may face public backlash. If citizens categorically abhor abandoning international partners in need, they should oppose Trump's signal about cutting off military aid for Ukraine. Conversely, if loyalty to external partners is more fickle, we should expect to observe the opposite effect when citizens are exposed to a potentially threatening signal of US withdrawal:
\paragraph{}
\textit{H2: The Trump Administration policy signal will increase support for raising the amount of national military aid to Ukraine.}

\subsection{The Cost of Alliance Cleavages}
The previous two sections have argued about the effects of cues from home and abroad on popular support for military aid. Despite a broad acknowledgment by the literature that policy signals from domestic \emph{and} foreign political actors compete for attention within a globalized information space, we still lack research that examines how competing messages \emph{across allied states} affect public support for war within alliances. 

At the level of international relations, competing signals within an alliance are a litmus test of whether alliance partners are willing and able to deliver on their commitments. To internal audiences, cohesion conveys assurance and a commitment to collective action, whereas cleavages convey discord and the threat of withdrawal.\autocite{jakobbarrett_how_2026} To external audiences, especially adversaries, cleavages within a security alliance communicate declining resolve and may even undermine the credibility of deterrence. The question remains how public support for military aid will fluctuate in the presence of competing signals, which communicate a loss of consensus among the coalition of the willing.  

Elite division over how best to end the war in Ukraine is likely to introduce uncertainty \autocite{berinsky_assuming_2007}. The literature on competing frames also suggests that when opinions are uncertain, exposure to competing frames can result in opinion change \autocite{matthes_diachronic_2012}. \textcite{chu_commanding_2021} find that negative rhetoric from political leaders erodes support for military alliances. However, they found that the effects dissolve in the presence of countervailing, positive rhetoric. This indicates that positive signals may be able to counteract negative signals, overcoming the potential negative effect of discord. 

Insofar as the public perceives NATO's commitment to Ukraine to be akin to supporting an ally in need, the literature on public opinion within alliances provides some insight into the penalties that leaders face when withdrawing. 
The literature indicates that foreign leaders pay a public opinion penalty not only for violating commitments to a shared ally \autocite{tomz_military_2021, chiba_careful_2015}, but also for breaking with the elite consensus. \autocite{chu_commanding_2021, reiter_manipulating_2024} If the public categorically opposes abandoning an embattled ally for moral reasons, they may also apply strong moral arguments for standing in solidarity with Ukraine. Competing signals from the domestic government and the Trump Administration about the future of military aid for Ukraine may therefore make the public \emph{less supportive} of abandoning an embattled ally. In other words, if we expect the Trump Administration to face public opinion penalties for withdrawing support from Ukraine, these penalties should not be neutralized when domestic governments oppose Trump's decision and argue that supporting Ukraine with military aid is \emph{the right decision}. Instead, following \textcite{chu_commanding_2021}, positive signals from the domestic government might counteract Trump's negative signals. 

However, the literature also suggests that leaders can minimize the public opinion penalties associated with withdrawing support from an ally. In addition to arguing that \emph{backing down} is the right decision \autocite{levendusky_when_2012}, leaders may find success in claiming that their country does not have an obligation to intervene on behalf of an ally \autocite{reiter_manipulating_2024}. In the domestic context, \textcite{reiter_manipulating_2024} demonstrated that the success of this strategy was conditional on whether political elites were unified behind the decision. Conflicting messages from elected leaders eliminated the persuasive value of the justification frame. Importantly, this study examined competing signals about alliance commitments in the domestic political sphere. We still lack empirical insight into how the public will react to competing signals from home and abroad.

Given conflicting theoretical expectations, I propose the third and final hypothesis. To the extent that discord between allies heightens uncertainty, we may expect the presence of negative rhetoric to erode support for solidarity. Conversely, assuming that the public perceives Ukraine to be an ally in need of support, we may expect positive signals from the domestic government to counteract Trump's negative signals, strengthening solidarity attitudes:
\paragraph{}
\textit{H3: In the presence of competing policy signals, support for the domestic government’s Ukraine aid policy will increase.}

\section{Research Design: Experiment}
To test these hypotheses, I fielded an original pre-registered survey experiment in the United Kingdom (n = 1528) and Germany (n = 1501). Participants were voting eligible, legal adults recruited from YouGov's survey panels. A final sample of 3029 was achieved after multiple attention checks screened out inattentive respondents.  

After the US, Germany and the UK are respectively the second and third largest donors of bilateral military aid to Ukraine. The Kiel Institute Ukraine Support Tracker calculates that US military aid to Ukraine amounted to 64.6 billion Euros sent between January 24th, 2022 and December 31, 2025. In the same time span, Germany has contributed 20 billion Euros worth of military aid, while the UK has contributed 14.3 billion. Analysis from the Kiel Institute suggests that any effort to replace US military support for Ukraine would first have to be shouldered by these two countries. \autocite{irto_ukraine_2025} Furthermore, compared to smaller European states, both the UK and Germany have broad security interests beyond their national borders, incentivizing both states to balance security threats within the larger scope of the European continent. 
Thus, understanding how the public of Europe's largest military powers responds when the US withdraws support for Ukraine is of both academic and societal importance. 

The commitment of both countries to providing Ukraine with such vast quantities of military aid is also made possible by the durability of a mainstream political consensus. In the United Kingdom, a cross-party consensus in support of Ukraine has been maintained in spite of inconsistent counter-narratives from Reform leader Nigel Farage. In Germany, the federal government has maintained a governing consensus in support of military aid for Ukraine, even after the 2025 election forced a change in the coalition makeup. Studying these two cases allows me to examine whether a consensus among governing parties is sufficient to sustain aggregate public support for Ukraine, even if AfD voters oppose the German government's policy based on their party's ideologically coherent message against military aid.  

Comparing these two countries with each other is informative because of their distinct strategic cultures, which are reflected in the mass public given stable trends in public opinion.\autocite{gravelle_structure_2017} Broadly speaking, the strategic cultures of the UK and Germany vis a vis the transatlantic relationship tend to fall into two different camps: Atlanticism and Europeanism.\autocite{dunne_when_2004, mader_national_2014, mader_no_2023,toettoe_identarian_2024} Survey evidence drawn from nationally representative samples in the UK has traditionally reflected a prevailing view that the United States is country's most valuable partner. \autocite{aspinall_uk_2023, aspinall_uk_2024} Despite recent rifts between the two countries and increasing public dissatisfaction with President Trump, a majority of the British public still believes the UK should either strengthen or maintain its current relationship with the US.\autocite{aspinall_uk_2025} The picture is quite different in Germany, where the public has traditionally been more supportive of strengthening European cooperation rather than cooperation with the United States. \autocite{mader_national_2014, mader_no_2023, gavras_nato_2020, mader_public_2020} German attitudes have soured on the US since the election of Donald Trump in 2024, with a majority expressing low confidence in the president to do the right thing in international affairs.\autocite{wike_us_2025} However, despite elite calls for a \textit{Zeitenwende}, evidence shows that the German public still wavers on whether or not their country should take up a stronger leadership role in Europe when it comes to security and defense.\autocite{mader_no_2023, poushter_german_2025} These contrasting strategic cultures provide theoretically motivated reasons to examine potentially heterogeneous responses to the same policy signals in Germany and the UK. 

\subsection{Vignette Treatments}
The experiment was fielded between February 13th to February 21st, 2025. The timing of this experiment took advantage of a rare chance to anticipate policy signals about the war in Ukraine before they became reality. Data collection was completed just days before tensions exploded during a White House meeting between President Trump, Vice President Vance, and the Ukrainian President Volodymyr Zelenskyy. One week after data collection, the Trump Administration officially halted the flow of military aid to Ukraine. The vignette treatments anticipated the actual policy signals eventually made by Trump and the national governments of both countries, therefore keeping hypothetical conditions close to reality as it was unfolding on the international stage. The vignettes announced an upcoming policy decision (hereafter referred to as a policy signal) from either the Trump Administration or the domestic government of the respondent's country about military aid for Ukraine. 

Both vignettes use a minimal frame around the policy signal, arguing that the decision to cut or increase military aid would be made \enquote{in order to put a stop to the fighting and create lasting peace and stability in Europe.} The minimal framing technique allows treatments to primarily manipulate information about the sender and the proposed policy, while holding the justification behind the policy constant. Following \textcite{levendusky_when_2012}, I chose the baseline justification that either increasing or cutting military aid is the \textquote{best way to end the war, put a stop to the fighting and create lasting peace and stability in Europe.} The tradeoff of this approach is that this study is unable to test how competing signals may interact with different types of justification frames. While the literature has identified many types of justification frames used by political leaders,\footnote{See \textcite{reiter_manipulating_2024} and \textcite{xu_reneging_2023}} a key strength of my approach is that the open-response data captures the different ways respondents independently reason about military aid in their own words. Comparing the control group with the three treatment groups (the latter of which received similar priming information about an impending decision about military aid) will thus allow this study to analyze whether treated respondents reason differently than non-treated respondents. 

The first round of randomization allocated respondents to four treatment conditions: a control condition, the domestic signal condition, the Trump signal condition, and the competing signal condition. The treatment texts are summarized in Table \ref{fig:treatment_table}. Through randomization, treatment groups were balanced on relevant background variables. For respondents in the competing signal condition, the order of the signals were randomized as well, to reduce potential ordering effects. Following the vignette texts, respondents in all conditions were asked whether they would support a decision by their government to increase military aid to Ukraine. This question, being the central dependent variable, was coded on a 5-point Likert scale ranging from strongly oppose to strongly support. 

  \begin{table}[H]
  \centering
  \footnotesize
  \begin{threeparttable}
  \caption{Vignette Treatments}
  \label{fig:treatment_table}
  \begin{tabular}{p{2.5cm}p{11.5cm}}
  \toprule
  \textbf{Treatment Group} & \textbf{Vignette Text} \\
  \midrule
  Control & Since the Russian invasion of Ukraine in 2022, [your country], the United States, and other NATO countries have sent multiple packages of military aid to Ukraine. \\
  \midrule
  Domestic Signal & Since the Russian invasion of Ukraine in 2022, [your country], the United States, and other NATO countries have sent multiple packages of military aid to Ukraine. However, in the coming
  months the (new*) government will likely reevaluate how best to end the war. Imagine that the government announces plans to increase military aid for Ukraine in order to put a stop to the fighting and create
   lasting peace and stability in Europe. \\
  \midrule
  Trump Signal & Since the Russian invasion of Ukraine in 2022, [your country], the United States, and other NATO countries have sent multiple packages of military aid to Ukraine. However, in his second term
  as President of the United States, Donald Trump will likely reevaluate how best to end the war. Imagine that President Trump announces plans to cut American military aid for Ukraine in order to put a stop to
   the fighting and create lasting peace and stability in Europe. \\
  \midrule
  Competing Signals & Since the Russian invasion of Ukraine in 2022, [your country], the United States, and other NATO countries have sent multiple packages of military aid to Ukraine. However, in his second
  term as President of the United States, Donald Trump will likely reevaluate how best to end the war. Imagine that President Trump announces plans to cut American military aid for Ukraine in order to put a
  stop to the fighting and create lasting peace and stability in Europe. Meanwhile, in the coming months [your government] will also likely reevaluate how best to end the war. Imagine that the government
  announces plans to increase military aid for Ukraine in order to put a stop to the fighting and create lasting peace and stability in Europe. \\
  \bottomrule
  \end{tabular}
  \begin{tablenotes}\small
  \item \textit{Note:} * For Germany, the new decision by the government explicitly referred to the new government, as the federal election was held just two days after data collection ended. The order of the
  signals in the Competing Signals condition was randomized.
  \end{tablenotes}
  \end{threeparttable}
  \end{table}

A number of secondary outcome measures were included in the latter half of the survey. Directly following the main outcome measure, respondents were asked to write a short, open-ended explanation of why they supported or opposed the domestic government's policy to increase military aid. Another secondary outcome measure asked respondents to indicate what they think other NATO countries should do regarding the amount of military aid they provide to Ukraine, with answer options including decrease by a lot, decrease by a little, not change the amount, increase by a little, and increase by a lot. 

A number of explanatory variables were measured prior to treatment to allow additional opportunities for exploratory statistical analyses, such as heterogeneous effects among various population subgroups. A handful of survey items tapped into future voting intention, past voting behavior, and ideology, all of which provide interesting insight into partisan-level effects of policy signals. The pre-treatment battery also included measures of foreign policy postures, which have been shown by the literature to be strong predictors of the types of specific foreign policies that individuals support. \autocite{wittkopf_foreign_1986, rathbun_taking_2016, gravelle_structure_2017,rathbun_dual_2020} Three individual survey items measured respondents' preferences to conduct foreign policy through international cooperation, military might, or isolationism. Additional explanatory variables in the latter half of the survey also allowed the design to explore secondary outcomes that may have been affected by the treatment.\footnote{A full version of the questionnaire can be found in the Appendix.}

\subsection{Modeling Approaches}
The three main hypotheses were tested by regressing the dependent variable, support for increasing in military aid to Ukraine, on the variable indicating treatment group assignment. Statistical analysis was performed in R, using ordinary least squares (OLS) regression models with robust standard errors. The main models (pooled and country-level) control for a pre-treatment measure of respondents’ opinions about military aid for Ukraine, which was measured prior to administering the vignette treatment. Including a pre-treatment measure that is closely related to the dependent variable has been shown to improve precision of the estimate by absorbing variance of the outcome not caused by treatment and increasing statistical power.\autocite{mutz_populationbased_2011, clifford_increasing_2021} A post-hoc test of the correlation between the quasi control and the post-treatment outcome variable supports this assumption.\footnote{See Table \ref{tab:correlation} in the Appendix} All variables were rescaled from 0-1 to ease the interpretation of effects.  

Additionally, I compared the estimates from the main model to a model including controls for age, gender, future voting intentions, and the three core foreign policy postures: isolationism, cooperative internationalism, and militant internationalism.\autocite{rathbun_dual_2020} In a randomized experiment, treatment assignment should be independent of all covariates by design. Adding controls to the model can improve precision by reducing noise, but even a simple difference-in-means already provides an unbiased estimate of the average treatment effect.\autocite{mutz_populationbased_2011} For all hypothesis tests, the estimates of both the main models and models with controls are presented, using standard inference criteria of p < .05.

To explore heterogeneous effects, the data is disaggregated and analyzed at the country-level, across political parties, and across respondents who self-reported strong alignment with one or multiple foreign policy postures. The same regression formula is used on the subsetted data, removing potential colliders from the list of controls to avoid multicollinearity and statistical bias. The following section presents the results of the pre-registered hypothesis tests.

\section{Results: Pre-Registered Experiment}
Figure \ref{fig:main_results} depicts the main results.

\begin{figure} [H]
    \centering
    \includegraphics[width=1\textwidth]{Figures/main_results_plot.png}
    \caption{Competing Cues Decrease Support for Military Aid}
    \label{fig:main_results}
     {\scriptsize {Note: Point estimates and confidence intervals are based on the model specifications described above. Bars represent the 95\% confidence interval. The coefficients estimate the effect of each signal when the control group is held as the reference category.}}
\end{figure}

Respondents who were treated with information about the domestic government's plan to increase military aid for Ukraine did not demonstrate a significant change in support for the policy. This finding aligns with the expectation of Hypothesis 1, that signals from the domestic government would not increase popular support for military aid. Although the public generally responds to policy signals about increased defense spending by swinging in the opposite direction,\autocite{wlezien_public_1995} the domestic signal about increased military aid for Ukraine did not generate strong opposition. On the contrary, it appears that thermostatic pressure manifests as a decline in support for military aid when voters receive policy signals from abroad. 

Information that the Trump Administration would announce plans to cut American military aid for Ukraine had a negative effect on popular support for military aid.\footnote{This effect was only found to be significant at the 90\% confidence level, but it remained robust to controls.} This runs counter to Hypothesis 2, which expected allied publics to punish the Trump Administration for backing out on Ukraine after NATO member states committed to supporting its defensive war against Russia. Within the context of this experiment, it seems that support for solidarity in Europe's two largest donor states did not increase in response to the US decision to break its commitments to Ukraine. 

Instead, abandoning an ally in need \emph{and} deviating from the alliance consensus actually appears to generate the opposite effect. When elites send competing signals about their willingness to support an embattled ally, the lack of alliance consensus results in a significant decrease in support for military aid. The decrease in support for aid amounts to -3\% and is significant at the 99\% confidence level.\footnote{The direction and significance of the effect is robust to controls. The power analysis in the Appendix indicates that both models are well powered at about 80\% to detect the negative effect of competing signals.} 

Taken together with the results of Hypothesis 2, it appears that the public does not oppose the Trump Administration's decision to abandon Ukraine. I reexamine this premise using approval of the domestic government and approval of Donald Trump as alternative dependent variables.\footnote{Regression results are included in Appendix Tables \ref{tab:gov_support_germany} and \ref{tab:gov_support_uk}} Surprisingly, none of the treatments affected approval of the domestic government or Trump in either country. This suggests that the public does not penalize Trump with lower approval ratings for leaving Ukraine in the lurch, as the literature on alliances and public support for war would suggest. \autocite{chiba_careful_2015, tomz_military_2021, levin_art_2022} Instead, as an external partner and aspiring NATO member,\autocite{noauthor_relations_2025} Ukraine may not benefit from the same type of loyalty mechanism that formal alliance members receive from allied audiences.\autocite{tomz_how_2023} 

\subsection{Country-Level Results}
When these effects are disaggregated at the national level in Figure \ref{fig:country_results}, the direction of effects is consistent with the main results, despite variation in significance across the two countries. 

\begin{figure} [H]
    \centering
    \includegraphics[width=0.8\textwidth]{Figures/country_plot.png}
    \caption{Country-Level Results}
    \label{fig:country_results}
    {\scriptsize {Note: Point estimates and confidence intervals are based on the model specifications described above. Bars represent the 95\% confidence interval. The coefficients estimate the effect of each signal when the control group is held as the reference category.}}
\end{figure}
%The pooled results indicate that citizens in NATO's two largest donor states do not rally behind Ukraine when the Trump Administration breaks with its commitments, nor when the domestic government's foreign policy is at odds with that of the US. 
\paragraph{}
As hypothesized, the domestic government's signal about increasing military aid for Ukraine does not significantly mobilize popular support in either Germany or the UK. Despite the potential presence of thermostatic pressure, exposure to the domestic government's policy signal about increasing weapons deliveries to Ukraine does not decrease \emph{or} increase public support for military aid. 

Consistent with the main results, information that the Trump Administration would cut military aid for Ukraine had no significant effect on German public support for military aid. However, the effect was negative and significant at the 90\% confidence level in the UK. \footnote{The direction of the effect was robust when controls were introduced in the model, and significance increased slightly to the 95\% confidence level} This appears to challenge a predominant interpretation in the literature that the public is morally and categorically opposed to withdrawing support from an embattled ally. Instead, it appears that the public in Europe's largest donor states became amenable to Trump's policy towards Ukraine, especially as the war entered its third year. Among German and British voters, mean support for increasing aid sits near the midpoint of the scale, indicating neither strong support nor strong opposition (see Figure \ref{fig:marginal_means} in the Appendix). 

A similar interpretation logically follows from the negative effects of competing policy signals. When the US signals that withdrawing support from Ukraine is the best way to end the war, it sows discord in the preexisting consensus of NATO's coalition of the willing. Neither the German or British government are able to mitigate the negative effect of consensus lost with a countervailing signal. As a result, support for military aid drops by 3.6\% in the UK sample (p < 0.05), mirroring the estimate detected in the pooled sample. While the significance washes down to the 90\% confidence level in Germany (and with the addition of controls in the UK), there is strong evidence for the negative direction of this effect in Europe's largest donor states. When discord within the alliance signals the permissiveness of withdrawing support from a NATO ally, the public follows suit. 

\subsection{Discussion: Do Competing Cues Create a Permissive Dissensus?}
%INTERPRETATION-Build up the Idea of a Permissive Dissensus
The main results indicate that support for military aid decreases in the UK and Germany when the domestic government's policy towards Ukraine clashes with that of the United States.\footnote{In the Appendix, I find that this effect is also significant for respondents with a militant internationalist posture towards international relations.} I argue that this particular effect is indicative of the power imbalance between the United States and junior partners in the NATO alliance. 

When a more powerful ally signals that stepping back from commitments to Ukraine is an acceptable and legitimate strategy, it permits voters to dissent from the mainstream political consensus and express opposition towards the practice of solidarity. This effect may be driven by prudence: voters might believe it is imprudent or impractical to continue sending military aid when the US pulls out, based on the size and importance of its contributions. 

However, the effect may also be driven by something other than rational and pragmatic considerations, which tend to require high cognitive burden. Competing signals are a sign of a broken consensus. In this experiment, the presence of competing signals conveyed to the public that Ukraine's three largest military aid donors disagreed about the best way forward as the war entered its third year. With no immediate end to the war in sight, disagreement between allied governments communicates a sense of uncertainty about whether military aid from Western powers will meaningfully alter the course of the conflict. Competing signals may create a permissive dissensus not only because voters are unsure if their government \emph{can} continue supporting Ukraine, but because voters are unsure if their government \emph{should}. 

Competing signals within alliances are not a rare or isolated phenomenon. Cohesion is notoriously difficult to ensure within alliances, especially as issue salience and the stakes of maintaining consensus increase. The Russia-Ukraine war is a specific case of a broken consensus. While a coalition of willing states informally agreed that standing in solidarity with Ukraine was both normatively and strategically necessary, the withdrawal of the coalition's most powerful state changed the stakes of solidarity as a foreign policy practice and increased the burden for the allies that remained.   

A burgeoning strand of research finds that US disengagement from its foreign policy commitments has spillover effects on foreign policy attitudes in Europe.\autocite{blankenship_threats_2024, jakobbarrett_how_2026} Despite a strong security threat from Russia increasing the risks associated with US defection, this experiment detected a tendency for disagreement between allies to cause public support for military aid to drop. But what types of arguments underlie that decline in support? 

The next sections dive into the individual arguments that respondents used to explicate their opinions on military aid. Below, I outline how the open-ended text responses were processed and coded. In this article, I present the main results that were derived from the manual coding procedure, and include further supplementary analyses (including natural language processing) in the Appendix. 

\section{Research Design: Coding Open-Ended Responses}
After reporting their support for increasing military aid to Ukraine, respondents were asked to explain their answer in a short, open-ended text response. The rich open-ended response data collected in this survey facilitates both a qualitative and quantitative analysis of the types of distinct justifications that voters use when explaining (in their own words) why they support or oppose the government's military aid policy. This choice also builds on a growing movement in political science and survey research to use open responses as a more granular source of data that elucidates, in voters’ own words, their considerations or reasoning behind political opinions.

First, open-ended responses were manually coded to allow for a more granular exploration of the different types of positive and negative arguments made by respondents. Following practices outlined by \citeauthor{krippendorff_content_2019}, the manual coding proceeded in three stages. All 3029 open responses were read and manually coded by the author\footnote{The author is fluent in both survey languages, raising no concern for misinterpretation in the initial coding phase.}. Although respondents were encouraged to explain their reasoning behind their opinion towards military aid, some respondents skipped the open-ended question entirely, wrote \textquote{I don't know}, or typed symbols and punctuation instead of words. Don't knows were preserved in the dataset as a unique category, while the remaining nonsensical or empty responses (n = 179) were discarded from the dataset and not coded. The remaining 2850 open-ended responses had an average length of 17 words, with a standard deviation of 15 words. 

These responses were coded by the author, first to distinguish from the open-ended text whether the respondent's attitude towards military aid was positive or negative. The substantive coding categories were created abductively by adding categories for new and distinct arguments as they appeared in the dataset.\autocite{krippendorff_content_2019} The coding unit was the entire open-ended response, which was coded for the presence or absence of unique arguments for or against military aid. The coding scheme allowed for mixed membership, meaning that a single open-ended response could  belong to multiple categories, depending on the number of unique arguments raised in the unit of text. Each open response was assigned a binary value (0,1) on each of the coding categories to reflect the presence of these distinct arguments.

The finalized coding scheme resulted in 60 unique argument categories, offering a more finely grained representation of the variance in responses. Inter-coder reliability checks were performed by a second coder fluent in both survey languages. The coder was provided the coding scheme and a randomly drawn sample of 50 responses, in order to test the replicability and logic of the coding categories. A reliability check was performed, comparing the agreement between the two sets of coded responses.\footnote{The average Krippendorff's \textit{alpha} was 0.84. Cases with low \textit{alpha} values were checked and re-coded in the final dataset to improve agreement between the coders.} Finally, the sixty coding categories were collapsed into nine major themes, which are discussed in the following section. 

To supplement the hand-coding and provide a meta-view of the different arguments made by respondents, the open-ended responses were also analyzed using natural language processing methods in R. These results can be found in the Appendix (see Tables \ref{fig:unigrams} and \ref{fig:bigrams}). The analysis of manually-coded open-response data follows as an additional and valuable exploration of the individual-level nuances of popular support for military aid. This exploratory analysis also points to important avenues for future research on how the public reasons about solidarity, obligations to allies, and competing signals within alliances.  

\section{Results: Text Analysis} 
Although a wide range of unique arguments were used to explain respondents' individual positions on military aid, the manual-coding also captured broader thematic clusters of justifications. These thematic clusters are discussed in Table \ref{tab:themecounts}.\footnote{See Table \ref{tab:argumentthemes} in the Appendix for a more fine-grained presentation of the themes and their component arguments, as well as the proportions of various themes across treatment groups and countries.} 

\begin{table}[H]
\centering
\small
\caption{Thematic Clusters}
\label{tab:themecounts}
\begin{tabular}{lr}
\toprule
\textbf{Themes} & \textbf{Occurrences} \\
\midrule
\cellcolor{gray!10}{\textit{Solidarity Arguments}} & \cellcolor{gray!10}{805} \\
\textit{Aid is Not Prudent} & 763 \\
\cellcolor{gray!10}{\textit{Domestic Cost \& Other Priorities}} & \cellcolor{gray!10}{737} \\
\textit{Strategic Arguments for Aid} & 716 \\
\cellcolor{gray!10}{\textit{Ambivalence}} & \cellcolor{gray!10}{277} \\
\textit{Not Our War} & 207 \\
\cellcolor{gray!10}{\textit{Anti-Trump / Pro-European Autonomy}} & \cellcolor{gray!10}{172} \\
\textit{Moral Opposition to War} & 90 \\
\cellcolor{gray!10}{\textit{Pro-Russia / Anti-Ukraine Sentiment}} & \cellcolor{gray!10}{89} \\
\bottomrule
\end{tabular}
\end{table}

The most prevalent theme (\textit{Solidarity Arguments}, n = 805) reflects responses that believe standing with Ukraine is \textquote{the right thing to do.} These arguments included appeals to democratic norms, moral obligations, and lessons from history. Responses were often explicit about the importance of \textit{sending a message to Putin} and other authoritarian leaders  that the invasion of another sovereign country would not be tolerated by European allies.

The second-most common theme in the dataset (\textit{Aid is Not Prudent}, n = 763) captures responses that question the prudence of sending military aid. This thematic cluster reflects skepticism about the effectiveness of military aid, with responses doubting whether continued arms transfers can meaningfully alter the course of the conflict or help Ukraine win. A considerable number of responses in this cluster argue that weapons only serve to prolong the war, that Europe cannot replace the quantity of military aid provided by the United States. Other arguments from this cluster stressed that sending military aid could further provoke Russian aggression against the UK or Germany and demanded a swift peace plan. Overall, this cluster conveys a feeling of fatigue towards sending military aid, given the conflict's long and seemingly intractable nature. 

A third prominent cluster (\textit{Domestic Cost \& Other Priorities}, n = 737) includes grievances rooted in the material tradeoffs of sending military aid and competing domestic or international priorities. Arguments in this category invoked economic burdens, unmet needs among the population at home, and frustration that other global crises currently receive insufficient attention. Responses often stressed that the government should assign higher priority to salient domestic issues affecting the everyday German or British citizen, including health care and education. A related but theoretically distinct cluster (\textit{Not Our War}, n = 207) reflects a sense of isolationism or disengagement, with responses emphasizing that the conflict falls outside German or British interests and that there is no meaningful obligation to intervene on Ukraine's behalf. This type of opposition is principled: respondents in these two clusters believe there are more pressing problems and that this is not their fight, regardless of whether or not aid would be effective. 

Strategic arguments in favor of sending military aid form a fifth cluster (\textit{Strategic Arguments for Aid}, n = 716), capturing responses that frame continued support in terms of containment and deterrence. Respondents argued that preventing Russian expansion and keeping the conflict from spreading beyond Ukraine's borders were pinnacle strategic priorities. Other responses within this cluster argue that increasing both the quantity of military aid and type of weaponry is precisely what is needed to bring the war to an end. 

Beyond these larger thematic clusters, three smaller but theoretically meaningful clusters emerge. A cohesive cluster of responses (\textit{Anti-Trump / Pro-European Autonomy}, n = 172) calls for greater European strategic autonomy in response to US disengagement, and criticizes the Trump administration for withdrawing support. A unique cluster of responses (\textit{Pro-Russia / Anti-Ukraine Sentiment}, n = 89) express sympathy toward Russia or outward hostility toward Ukraine and Ukrainian people, representing the most extreme oppositional arguments in the data. Another cluster (\textit{Moral Opposition to War}, n = 90) expressed clear moral opposition to military aid, given that it constitutes financing warfare and by extension, the arms industry. These moral arguments stress that Europe should not be financing Ukraine's war as a matter of principle, rather than prudence.

Finally, a non-trivial share of responses (\textit{Ambivalence}, n = 277) express ambivalence, a lack of knowledge, or that the situation is too complex to take a clear position.

\subsection{Treatment Effects on Argument Probability}
Coding arguments into larger thematic clusters allows for additional statistical analysis into the marginal probability of invoking a particular type of argument, depending on treatment. I used a set of OLS regression models using robust standard errors, mirroring the specifications of the estimator used to calculate the main treatment effects in the experiment. The model with controls uses the same set of covariates (age, gender, vote intention, foreign policy postures, and pre-treatment attitudes towards aid). Each thematic cluster was treated as a dependent variable, so that the estimate can be understood as the marginal effect of treatment on the probability of invoking an argument within that given theme. 

I modeled the effect of receiving either the domestic policy signal, the Trump policy signal, or competing policy signals (relative to the control group). I ran the same models (with and without controls) at the country-level to check for heterogeneous effects. 

%Figure \ref{fig:treated_theme} depicts the marginal effect of being treated on the probability of invoking a theme. Two significant effects are immediately apparent: being treated makes respondents not only more likely to argue that the conflict in Ukraine is \textquote{not our war}, but also less likely to invoke solidarity arguments. However, the confidence intervals are quite large, indicating that the estimate itself is not very precise. 

Figure \ref{fig:treatment_theme} displays the effects of each treatment relative to the control group, in order to explore whether certain policy signals made respondents more likely to invoke different arguments. Competing signals make respondents 6.7\% less likely to argue for the importance of standing in solidarity with Ukraine, an effect which is significant at the 99.9\% confidence level. Respondents in the competing signals treatment group are also 3.6\% more likely to argue that the war in Ukraine is removed from German or British interests (p < 0.01).\footnote{Regression results for these two argument clusters can be found in Tables \ref{tab:notourwar} and \ref{tab:solidarity} in the Appendix.}

Interestingly, respondents in the domestic policy signal condition also became less likely to invoke solidarity arguments. This finding provides additional evidence that the domestic policy signal is ineffective at rallying the public behind supporting Ukraine, neither on its own nor in juxtaposition with a countervailing signal from the Trump Administration. Meanwhile, Trump's policy signal did not have significant effects on any argument themes. 

\begin{figure} [H]
    \centering
    \includegraphics[width=1\linewidth]{Figures/treatment_theme_plot.png}
    \caption{Competing Signals Make Respondents Less Likely to Argue for Solidarity}
    {\scriptsize {Note: Point estimates and confidence intervals are based on the model specifications described above. Bars represent the 95\% confidence interval. The coefficients estimate the effect of being treated with any signal when the control group is held as the reference category.}}
    \label{fig:treatment_theme}
\end{figure}

When these effects are calculated at the country-level, an interesting pattern emerges. Figure \ref{fig:treatment_theme_country} shows a similar pattern to the one detected in the aggregated sample. However, the models also indicate that competing policy signals prompt different arguments in Germany and the UK. 

\begin{figure} [H]
    \centering
    \includegraphics[width=1\linewidth]{Figures/treatment_theme_country_plot.png}
    \caption{Competing Signals Prompt Different Arguments in Germany and the UK}
    \label{fig:treatment_theme_country}
\end{figure}

We can see that competing signals increase the likelihood of making \textquote{not our war} arguments by 4.2\% in Germany (p < 0.05), but the estimate does not reach significance in the UK.\footnote{See Table \ref{tab:notourwar_de} in the Appendix} The same treatment significantly decreases the likelihood of invoking solidarity arguments in the UK by 9.5\% (p < 0.01), but the estimate does not reach significance in Germany.\footnote{See Table \ref{tab:solidarity_uk} in the Appendix} Nonetheless, the direction of effects appears consistent with the estimates calculated in the aggregated sample. 

Interestingly, the domestic signal also increases the likelihood that British respondents express ambivalence towards the issue of military aid by 6.7\% (p < 0.01).\footnote{See Table \ref{tab:ambivalence_uk} in the Appendix} This is an interesting effect in itself. In the absence of competing policy messages, the domestic government's proposal to increase military aid for Ukraine prompts respondents to become more uncertain, relative to those who were not treated. 

Taken together, the text analysis provides useful exploratory evidence into the considerations that respondents make when prompted with various policy signals. Competing signals from home and abroad decrease support for military aid across both countries, and decrease voters' willingness to make arguments in favor of standing in solidarity with Ukraine. This is consistent with the logic of the permissive dissensus. When the US signals that it is no longer in American interests to support Ukraine, it prompts citizens in allied countries to express a similar opinion.  

\section{Conclusion}
This study provides evidence that support for military aid declines in the presence of competing signals from home and abroad. I argue that these results indicate the presence of a permissive dissensus. 

At the time of fielding this experiment, the United States was the largest donor of military aid to Ukraine, and it remains the most powerful member of NATO. When the Trump Administration announced that backing away from the war in Ukraine was an acceptable and legitimate strategy, it created an opportunity for citizens in allied states to dissent from the mainstream political consensus and voice opposition towards the practice of solidarity. This is apparent from the negative and significant effect of competing signals, which directly juxtaposes two allied states' approaches to solidarity and military aid. 

Through manual coding of open-response data, this study sheds additional light on the tradeoffs and risks that voters associate with solidarity for allies in need. Opponents of military aid were most likely to question the prudence of weapons deliveries and express dissatisfaction with the material and domestic tradeoffs of sending aid. OLS regression models estimated that competing signals do in fact offer a permissive environment for voters to express grievances about the practice of solidarity. Allied disagreement over military aid significantly reduced the likelihood that voters would invoke solidarity arguments, and increased the likelihood of arguing that the UK and Germany are categorically not obligated to fight Ukraine's war. 

%This study also makes a key contribution to the literature on alliances and public support for war. Previously, the prevailing wisdom in the literature was that the public punishes leaders who break alliance commitments, especially if it means abandoning an ally that has been unfairly attacked. \autocite{chiba_careful_2015, tomz_military_2021, levin_art_2022} The results of this experiment indicate that the public does not hold allied leaders accountable for backing away from an ally in need. Instead, public opinion tends to follows suit, particularly when the domestic government's policy contrasts with that of a more powerful allied state. In sum, leaders may indeed face public opinion penalties if they abandon an ally in need. Previous research has provided evidence that domestic leaders are punished by their own constituents for reneging on alliance agreements,\autocite{tomz_military_2021, levin_art_2022, xu_reneging_2023} indicating that domestic leaders might be more vulnerable to such penalties than their foreign counterparts. However, when I examined this mechanism in a highly realistic experimental setting, I found that the desire of public audiences to hold foreign leaders accountable wanes when war fatigue and thermostatic pressures sink in. 

The findings of this experiment reflect a larger phenomenon about the fickle nature of cohesion in asymmetric alliances. Competing signals between allied states reveals a fractured consensus, which has measurable consequences for public opinion. Competing signals communicate to the public that there is visible disagreement among Ukraine's principal backers, which permits citizens to doubt whether their governments should sustain support amid such uncertainty. The Russia-Ukraine war illustrates this dynamic with particular clarity. What began as an informal coalition united around the strategic and normative case for solidarity was fundamentally destabilized when its most powerful member withdrew, raising the costs of commitment and casting doubt over whether European military assistance could meaningfully alter the course of the conflict.

Importantly, the results of this experiment also reflect specific alliance power dynamics. The United States is uniquely positioned to generate a permissive dissensus, given that its decision to halt military aid for Ukraine affected the security risks and financial costs associated with filling the gap. These dynamics will continue to challenge how the remaining coalition of willing states sustain public support for solidarity in such a protracted conflict. Moreover, these particular scope conditions may travel to other cases in which a dominant or more powerful ally defects from commitments, creating a window for junior allies to either step up or step back. Future research should examine whether the permissive dissensus requires hegemonic power to take hold, or whether the effect is replicated when a weaker ally withdraws from alliance commitments. The permissive dissensus may well have similar implications for other forms of international solidarity, such as humanitarian aid, disaster relief, and climate adaptation, which are equally important to study. 

Public audiences are clearly sensitive to competing cues from home and abroad. This study offers a first insight into how exposure to alliance cleavages prompts citizens to oppose the practice of solidarity and voice grievances against the status quo. Public support is shaped by the interpretive signals that allied governments provide; when those signals diverge, an informal consensus among allies risks losing the public legitimacy it needs to persist. This suggests that the sustainability of consensus among allies depends not only on material contributions, but on the coherence of political narratives. Alliance cohesion, in this sense, is only as strong as the signals sent by its most powerful members. 

%TC:ignore
\clearpage
\printbibliography
\clearpage
\input{3-appendix}
\clearpage
%TC:endignore
\end{document}